\chapter{Introdución}
\label{chap:introducion}

\lettrine{E}n los últimos tiempos, la robótica autónoma ha experimentado un gran avance gracias 
al desarrollo de nuevas técnicas de aprendizaje automático, las cuales permiten a los robots 
adaptarse y operar en entornos más complejos y cambiantes que las aproximaciones más tradicionales.

El método que trataremos a continuación busca conseguir que los robots sean capaces de aprender 
de forma autónoma modelos del mundo y objetivos en base a experiencias propias. 
En este contexto, los modelos de mundo permiten al agente, que se encuentra en un entorno y en cierto estado 
en el instante actual, predecir el siguiente estado tras haber realizado 
una acción; mientras que los modelos de utilidad buscan discernir el nivel de eficiencia que 
tendrán dichas acciones para alcanzar cierta meta. La combinación de ambos da lugar a un sistema 
deliberativo de toma de decisiones capacitado para completar estas tareas.

En este trabajo nos centraremos en la creación de este sistema, basado en modelos de mundo 
y utilidad, aplicado a un robot cuadrupedo Unitree Go2\cite{UnitreeGo2Web}, para completar una tarea dada al 
poder seleccionar deliberadamente la mejor acción de un conjunto de acciones para un estado dado. Este deberá ser capaz de identificar 
objetivos de forma autónoma (en este caso, aprender a acercarse o alejarse de un objetivo) y cumplirlos en un 
entorno real (reducir la distancia a dichos objetos seleccionando la acción más adecuada para ello).

\section{Motivación}
\label{sec:motivacion}

La motivación principal de este trabajo nace de las limitaciones que presentan los enfoques tradicionales 
de la robótica basada en control clásico, en los cuales, características clave como la función de recompensa, los objetivos o 
los modelos del entorno son definidas por los programadores. El problema de este tipo de aproximaciones 
es que, aunque funcionan bien en escenarios controlados, tienen poca capacidad de adaptación y generalización 
ante nuevos escenarios o metas distintas.

Ante esta situación, el aprendizaje de un modelo deliberativo, es decir, la combinación de modelos de mundo 
y utilidad, supone una excelente alternativa, permitiendo al agente ser capaz de descubrir objetivos y 
aspectos relevantes del entorno.
Trabajos previos, algunos de los cuales comentaremos más adelante (sección \ref{sec:estado_arte}, página \pageref{sec:estado_arte}), 
han demostrado que el uso de modelos de mundo mejora en gran medida la robustez y la eficacia de los sistemas 
en entornos complejos y con información parcialmente incompleta. Asimismo, los modelos de utilidad permiten, la independencia y reutilización 
de los modelos de mundo y el descubrimiento de múltiples metas, cada una evaluada con su propio criterio.

Desde el punto de vista académico, este trabajo permite una aproximación a múltiples tecnologías relacionadas 
con la informática, como el aprendizaje automático, la robótica, la simulación y el desarrollo de software. 
Además, la realización de este en un entorno real con el robot y las cámaras Optitrack, así como en uno simulado, 
permite apreciar las diferencias en ambos tipos de desarrollo y las dificultades para trasladar los conocimientos 
adquiridos entre ambos entornos. 


\section{Objetivos}
\label{sec:objetivos}

Lo que se pretende lograr con este trabajo es desarrollar un sistema aprendizaje autónomo de un modelo deliberativo 
de toma de decisiones, basado en modelos de mundo y de utilidad, para un robot cuadrupedo, que le permita encontrar 
y aprender a alcanzar una serie de objetivos descubiertos previamente.

Para hacer realidad esta meta se definieron una serie de objetivos más concretos:

\begin{itemize}
  \item Aprender autónomamente un modelo de mundo, basado en el dominio en el que opere el robot, que permita predecir 
  las percepciones futuras en base a las percepciones actuales y a la acción que se decida tomar.
  \item Aprender autónomamente un modelo de utilidad que evalúe la utilidad de las percepciones del robot con respecto 
  al objetivo que se desea alcanzar. Será necesario crear un modelo de utilidad independiente para cada uno.
  \item Implementar este aprendizaje autónomo en un entorno simulado, en este caso MuJoCo\cite{MujocoWeb}, para facilitar 
  el análisis de estos modelos.
  \item Transferir estos modelos de la simulación a un entorno real con un robot cuadrupedo Go2 de 
  la marca Unitree\cite{UnitreeWeb}, utilizando un sistema de captura de movimiento Optitrack\cite{OptitrackWeb} para 
  obtener la posición en el entorno.
  \item Validar, comparar y analizar los resultados obtenidos tanto en el entorno simulado como en el entorno real.
  \item Validar y analizar la independencia entre modelos de mundo y de utilidad.
\end{itemize}



